\section*{Author Appendix I, Algorithm 1}
Integer factorization assuming the truncations of $(1+k^2x)^{i/k}$ being easy for each $i$. Input: A composite positive integer $n$. Output: A non-trivial factor of $n$.
\begin{enumerate}
\item Define $N(d,k):=k^{(k-2)d}(dk)!/(d!)^k$.
\item $m\leftarrow k$.
\item \textbf{while true do}
\item Compute $\gcd(N(m,k),n)$.
\item \textbf{if} $\gcd(N(m,k),n)=1$ \textbf{then}
\item $m\leftarrow mk$.
\item \textbf{else if} $\gcd(N(m,k),n)=n$ \textbf{then}
\item $t\leftarrow m$. This $m$ is the desired $t$.
\item \textbf{break}
\item \textbf{else}
\item \textbf{return} $\gcd(N(m,k),n)$. Here $\gcd(N(m,k),n)$ is a non-trivial factor of $n$.
\item \textbf{end if}
\item \textbf{end while}. At this step, all the primes dividing $n$ are in the interval $[t+1,tk]$.
\item $a\leftarrow1$.
\item $b\leftarrow t$.
\item \textbf{while true do}
\item \textbf{if} $b-a\le1$ \textbf{then}
\item $s\leftarrow a$. This $a$ is the desired $s$.
\item \textbf{break}
\item \textbf{end if}
\item $c\leftarrow\lceil(a+b)/2\rceil$.
\item Compute $\gcd(N(c,k),n)$.
\item \textbf{if} $\gcd(N(c,k),n)=1$ \textbf{then}
\item $a\leftarrow c$.
\item \textbf{else if} $\gcd(N(c,k),n)=n$ \textbf{then}
\item $b\leftarrow c$.
\item \textbf{else}
\item \textbf{return} $\gcd(N(c,k),n)$. Here $\gcd(N(c,k),n)$ is a non-trivial factor of $n$.
\item \textbf{end if}
\item \textbf{end while}. At this step, all the primes dividing $n$ are in the interval $[sk+1,(s+1)k]$.
\item \textbf{for} $i=sk+1$ \textbf{to} $(s+1)k$ \textbf{do}
\item \textbf{if} $i$ divides $n$ \textbf{then}
\item \textbf{return} $i$. Here $i$ is a non-trivial factor of $n$.
\item \textbf{end if}
\item \textbf{end for}
\end{enumerate}
